\section{The Exoteric Path}

\begin{quotex}
If our contemporaries as a whole could see what it is that is guiding them and where they are really going, the modern world would at once cease to exist as such, for the rectification could not fail to come about through that very circumstance. \flright{\textsc{Rene Guenon}}

A day will come, one day in the unending succession of days, when beings, beings who are now latent in our thoughts and hidden in our loins, shall stand upon this earth as one stands upon a footstool, and shall laugh and reach out their hands amidst the stars. \flright{\textsc{H G Wells}}

\end{quotex}
\textbf{Rene Guenon} appeared in the world process at the time of the spiritual emptiness of the West. With all of history and geography pretty much known, it had become impossible to ignore the vast differences in spiritual experiences in the world. Guenon was able to sort them out and to reveal their deep similarity in terms of social organization, inner experience, and metaphysical principles. Yet his voluminous works are just a beginning. The logical consequences of his writings still need to be explicated. He himself suggested several areas for further research in those writings.

Although we don't always make them explicit, we conform to the core of Guenon's principles. That is not a slavish following, a ``Guenonian scholasticism" as Evola called it. Evola himself, although calling Guenon the Master of the Twentieth Century, challenged Guenon on certain points. Sometimes he was mistaken but on other times he made key contributions. Furthermore, we need to distinguish between metaphysical principles which are apodictic and personal judgments which are at best informed opinions.

\paragraph{Guiding Principles}
\begin{quotex}
If worship of providence were forbidden, the natural consequence would be their fall, for a nation of fatalists or casualists or atheists never existed in the world, and we saw that all the nations of the world believe in a provident divinity, and have embraced only four principle religions: paganism, Judaism, Christianity, Islam. \flright{\textsc{Giambattista Vico}}

\end{quotex}
We are not out to ``prove" any Guenonian principles; for that, his own works need to be consulted. Rather, we accept them, even if provisionally, and draw out their consequences as if they are true.

First of all, the modern world is an anomaly in the world, as no civilization has ever been based on its principles, or, better said, anti-principles. Instead, they have all been based on what he calls Tradition, whose basic principles have been described throughout Gornahoor. Tradition means ``handed down", not a nostalgic return to some idealized past. Hence, it is not necessarily Luddite, as the example of the City of the Sun\footnote{\url{https://www.gornahoor.net/?p=518}} shows.

Traditional societies are hierarchical with specific functional areas and daily life was centered around the rites and rituals of an exoteric religious tradition. Beyond that, there is an esoteric tradition related to, but different from, the exoteric tradition. The esoteric tradition comprises various practices and metaphysical teachings that lead to higher states of being.

\paragraph{Exoteric and Esoteric}
The relationship between the exoteric and the esoteric is expressed in this principle: every exoteric religious or theological statement is equivalent to an esoteric metaphysical statement. That is why we don't get involved in theological disputes since they can't be resolved at that level without appealing to faith or authority. Esoterism, on the other hand, results in gnosis, or certain knowledge. In other words, what the believer believes, the knower knows.

For example, the religious idea of monotheism is the metaphysical idea of the Absolute, i.e., there is one foundational principle of the world. Polytheism, hence, means competing principles with no underlying unity. Atheism is the denial of an underlying principle so everything that happens is the result of determinism or chance events. As Vico pointed out, no society based on atheism, fatalism, or chance events could survive for long. The modern mind considers atheism as simply another defensible opinion, but for the traditional mind the atheist is an enemy of the people and a danger to the state. He would have no more right to free speech than a quack physician who prescribes harmful and ineffective treatments.

Guenon's next point is that the esoteric component has been lost in the West. We accept that, but not absolutely since it is an opinion not a metaphysical truth. Nevertheless, it is pointless to criticize bishops, popes, and so on, for that reason. The task, for us, is to delineate the features of a Western esoteric Tradition to prepare the way for its return.

A further point is that the esoteric should not be seen in opposition to the exoteric. Rather, those capable of and interested in pursuing and esoteric path should do it in the context of an exoteric path. These would be, at the risk of over-simplification, one of the many paganisms or one of the Abrahamic religions as Vico noted.  The normal choice, however, is to choose the exoteric religion of one's land. Hence, Guenon did not fully embrace Islam as a lifestyle until he migrated to Egypt.

The notion of picking and choosing and shopping around for one's religion is absurd from the Traditional point of view. Religious pluralism is a liberal and modern idea. The opposite Traditional view is called ``Integralism", i.e., the nation as an organic unity. Evola's solution of avoiding any commitment to a Tradition was an unfortunate choice and serves as a bad example.

A writer can praise the virtues of loyalty and fidelity in the abstract, but they must be tested in practice. Hence, a commitment to an exoteric religion can be a serious challenge to one's loyalty and fidelity. However, we accept all its dogmas and rites as an act of will. (Faith is an act of intellect commanded by the will.) The point is that we are not out to reform the exoteric religion but rather to bring to light the latent esoterism in it.

We often hear the complaint that the exoteric teachings are not good enough or that there is not esoteric teaching, and on and on. That, for us, is simply too passive; if you are not finding anything better, then you probably don't deserve anything better. A more active approach is to stifle the complaining while mining the hidden gold nuggets of esoterism.

As a reminder, anyone who is thriving under the exoteric religion does not need anything here. Esoterism is purely optional and is intended for those to whom the exoteric religion, as conventionally expressed, lacks all power and meaning.

\paragraph{Paganism}
As we have seen, pagans considered that there was a semi-divine founder of their tradition who established the spiritual, religious, and political practices that dominated their everyday life. By maintaining the worship of the ancestors, and appeasing or propitiating the gods, providence would continue to guide the people. Hence, we see that the Indo-Aryan Brahmans could trace their lineage back to 21 rishis, each Greek city had its own divine founder, and Rome was founded by a few noble families.

Hence, any attempts to revive Western paganism are counter-traditional, unless a demigod arrives to start a new family. It is not possible to just jump into an existing stream. The other issue is that in a pagan society, the bulk of the population would be serfs, slaves, or shudras. No revivalist cares to mention that.



\flrightit{Posted on 2015-02-04 by Cologero }

\begin{center}* * *\end{center}

\begin{footnotesize}\begin{sffamily}



\texttt{Lu Dongbin on 2015-02-05 at 09:03 said: }

How does traditional China, particularly the Tang Dynasty, stand in relation to your notion that ``picking and choosing" ones religion is a modern and liberal idea? During the Tang people were free to choose from a variety of religions which coexisted with each other such as Daoism, Buddhism, Nestorian Christianity, Manichaeism, and Islam (Confucianism too if you consider it a religious tradition prior to the metaphysical reform of so-called neo-Confucianism.)

While some of these died out, many continued to be present on Chinese soil and choices available until the Red takeover.


\hfill

\texttt{Cologero on 2015-02-05 at 09:17 said: }

First of all, Lu Dongbin, we are focusing on the West and the recovery of Tradition therein. I am not familiar enough with Chinese spirituality, but I believe that it, too, at least on one point was founded on ancestor worship.

Second of all, pluralism exists in an Empire; we specifically referred to the ``nation". For example, the Roman Empire comprised many nations and religious forms, but the Evola essays we recently translated referred to the earliest periods of Rome.

According to Guenon, Confucianism was the exoteric form and Daoism the esoteric form of the Chinese religion. Buddhism is a special case, since it was at its inception primarily a metaphysic and a path of self-development. As such, it could be absorbed into other traditions, as happened in China and Japan. In India, its metaphysical teaching become incorporated into Advaita Vedanta so it barely survives as a separate religion.

Christianity and Islam were brought in by foreigners and have not historically defined the Chinese consciousness. (Maybe that will change in the future?) There will always be pockets of dissent even in an integral society.


\hfill

\texttt{Lu Dongbin on 2015-02-05 at 10:10 said: }

I am not Chinese myself but rather a Westerner interested in Chinese traditions. I realize your focus is on the West and the recovery of Tradition (Christianity) here, but how should one view the new traditions that have planted their feet in the West in recent years such as Buddhism, etc. In the scenario Gornahoor's project were to succeed, would these new traditions have to uprooted, would they simply disappear, or would a situation akin to the Tang Dynasty I mentioned earlier have to be the way to approach it?

Also, just a few notes, while it is the case that China was an empire, it was a bit different in this case than Rome was due to the overwhelming dominance of the Han ethnic group, whereas of course in Rome you had Latins, Britons, Greeks, Syrians, Egyptians, etc. in reasonable numbers. In this sense the empire was nearly equivalent to a nation despite including other minority ethnic groups.

Secondly, the Buddhism that arrived in China wasn't the pure metaphysic of earliest Buddhism but rather the religious tradition of Mahayana, complete with its own pantheon, cosmology, eschatology, etc. I am not sure if we can say it was absorbed in a manner similar to Advaita because while it left a massive influence on Daoism (which in turn influenced it and helped create Ch'an/Zen), it also competed with the other traditions at various points in history and remains a separate tradition. Many Chinese tended to approach Daoism, Confucianism, and Buddhism as a harmonious unity, the so-called idea of ``san jiao he yi" and would often learn from all three.

An interesting discussion about the Chinese traditions in relation to Guenonian Traditionalism can be found here, particularly the exchanges between a user named ``Justice\&Mercy" and Charles Upton:

\url{http://www.sophiaperennis.com/discussion-forums/neglected-traditions/what-about-neglected-traditions/}


\hfill

\texttt{Cologero on 2015-02-05 at 11:52 said: }

In a spiritual void, who can say what will prevail? I've referenced Guenon's take\footnote{\url{http://www.gornahoor.net/?p=35}} on it several times. For example, Michel Houellebecq's new novel (in French, but English translation expected in the fall) shows how the Islamization of France may occur. He refers to Guenon in that book, by the way, so Guenon will be going mainstream in intellectual circles. Be prepared.

I accept Evola's list of Western traditions: Brahmanic India, pagan Rome/Greece, and Medieval Europe. Those can be reincorporated into a larger whole, perhaps along the lines laid out in Guido De Giorgio's ``Vedantized Christianity". I'll hold off on that discussion until the second part.

Buddhism has become quite decadent in the West. Many people I know call themselves ``Buddhists" because it sounds cooler than atheist, but they hardly believe that all life is suffering. Nor do they accept the sexual restraints that are part of their vows of initiation. The Dalai Lama has turned out to be a rather weak fellow, capitulating to Western values at every opportunity.

What is happening in the USA is a New Age amalgam based around the Course in Miracles, the Power of Now, and the Law of Attraction. Speaking of Upton, he wrote a long book (The System of the Antichrist) on several New Age movements from a Traditional perspective. He gave them more attention than they deserved and actually made them some appealing and even traditional in several aspects. Since I've been looking at that book lately, I checked Amazon. Miracles is \#5 on Amazon's ``Christian Living" category.

I don't accept Justice\&Mercy's approach. Just because something happened in the past does not make it ``Traditional". We are interested in the metaphysical not the historical; the latter is used to illustrate the former. And we can hardly be concerned about the religious practices of Tom, Dick, and Harry, but the focus should instead be on the highest exponents of a given tradition.

Furthermore, his partisanship is not really helpful. He mentions Guenon's Great Triad book with high praise. But Guenon's point is that the teaching found in Daoism is also found in other traditions.


\hfill

\texttt{Lu Dongbin on 2015-02-05 at 15:26 said: }

Well I certainly agree with your comments on Western Buddhism–it is largely, though not entirely, a joke full of materialists, atheists, those who conflate modern politically correct leftism with ``the dharma", and those who try to psychologize Buddhism (especially Zen) as a means to cope with daily problems or find peace within daily life (not bad aims, but not enough) rather than awaken to the Absolute and intuitively uncover the answer to that grand question, ``Who am I?" Then again, not many modern Christians follow the ethical standards of the ancients or are searching for a means to begin the process of theosis either.

I should point out that I also don't entirely agree with that fellow ``Justice\&Mercy" mentioned, but I think he may have been on to something in that the Traditionalism of Guenon and others is primarily from the viewpoint of Hinduism, Christianity, and Islam. It seems many of the categories or ideas Guenon and others held cannot be consistently applied within the Chinese sphere with their multiplicity of traditions, syncretism, etc. (not limited to Tom, Dick, and Harry from my studies) but who nonetheless produced many sages over the centuries and maintained highly metaphysical aspects of their culture until relatively recently.

On a personal level, I was born without a tradition and raised in a secular household and have been searching for where to focus my efforts these past 5 years or more. While I continue to appreciate Guenon, Evola, and the Traditionalists and agree with most of their views, in finding which tradition to devote my life to I don't start from the perennialist point of view and work back but rather try to understand the particularities of each tradition and discern which are true and possess the most efficacy/potency to lead me to gnosis and salvation. This may have put me in the ``pick and choose" category until just recently, but the unique conditions of my time, nation, and birth forced me into such a category.


\hfill

\texttt{ja on 2015-02-05 at 12:23 said: }

Couldn't Aleister Crowley be seen as a reviver of the pagan stream, through his contacts with Aiwass the god of ancient Sumeria ? Playing devil's advocate here….


\hfill

\texttt{Cologero on 2015-02-05 at 22:23 said: }

Sure, Ja, and I know Thelema has many followers today. And there are other dictated teachings such as the Seth books or Alice Bailey. The difference, I think, is that some spirit has revealed something. But I don't believe that any one of them is a spiritual ancestor in any sense, neither by birth nor adoption.


\hfill

\texttt{SydneyTrads on 2015-02-07 at 22:38 said: }

We would be grateful if a pinpoint reference could be provided for Giambattista Vico's indented quotation.


\hfill

\texttt{Cologero on 2015-02-07 at 23:28 said: }

The quote is from Vico's New Science, paragraph [602].

\begin{quotex}
If worship of providence were forbidden, the natural consequence would be their fall, for a nation of fatalists or casualists or atheists never existed in the world, and we saw that all the nations of the world believe in a provident divinity, and have embraced only four principle religions: paganism, Judaism, Christianity, Islam. 

\end{quotex}

\hfill

\texttt{KSS on 2015-02-10 at 12:23 said: }

``The normal choice, however, is to choose the exoteric religion of one's land." 

What if the exoteric religion of one's land is cultural marxism? What if the values of your country's ``christian church" are identical with those of the new left radicals?


\hfill

\texttt{Cologero on 2015-02-10 at 13:02 said: }

Please read the follow up post to this one, which goes into more detail.

First of all, ``cultural marxism" is not a valid tradition.

Then, as Guenon explains, ``its present representatives seem to be unaware". The Tradition is one thing, its management is something else. We are interested in preserving the former, not necessarily the latter.

Also Guenon explicitly writes: ``In this case it would be better, although not absolutely necessary …", so if you have a better choice, then go for it.

The bottom line, KSS, is that if you see yourself — much like a typical citizen of the model world — as a passive consumer waiting for George to create a suitable Tradition ready-made for you, you will wait in vain. What is necessary are active and creative agents.


\hfill

\texttt{Cologero on 2015-02-16 at 09:43 said: }

Let me tell you some personal stories, Nyklot.

Yesterday, I was walking along the beach, checking out the girls every now and then, when I was drawn to the ocean. So I sat down and did a meditation. I was struck by its size and power and looked out onto the horizon. I wondered about the character of the first men who determined to sail out onto that horizon. I tried to feel the spirit of the ocean, since everything is alive. Eventually, in the depths, I experienced the presence of Poseidon and considered what had been lost or gained. But Poseidon was himself subject to destiny. Deeper still, I felt the living presence of the God. So nothing was lost, but the overcoming of destiny was gained. Unfortunately, in our time ``God" is out there, not as a living presence, so the ocean is just the movement of matter.

A second story involves a trip I took with my son several years ago to Sicily to explore our roots. My family had lived in Salemi for 5 or 6 hundred years since the arrival of the first Salvo from Tuscany. Now Salemi was the site of an old castle used by the Normans when they re-established the Holy Roman Empire from Sicily. High on a hill, one could see everything for several kilometers. Not quite sure what I was looking for, I checked the walls, both inside and out, looking for Runic inscriptions. There were none … that time had long past, even one thousand years ago.

There was a much larger castle in Palermo, where the situation was clarified. The castle had a beautiful chapel, the Palatine Chapel. I was awestruck, as I mentally put myself as one of the knights in the presence of that chapel. That was what was flowing in their blood, not some lifeless paganism.

So, Nyklot, you can dismiss a thousand years of your ancestors as somehow being ignorant, or deluded, or whatever unspoken insult you want to hurl on them. But if you truly want to honour them, you should honour their choice. At one time there was a Europe that was spiritually united, while respecting the multiple nationalities.

Furthermore, Nyklot, you may want to consider why, or by whom, medieval Christianity has been under siege for centuries. In that so-called clash of civilizations (which is ongoing), where do you take your stand? An effete and deracinated paganism, suitable only for fanboys, is no threat to anyone. There are no ``cartoons" satirizing pagans.

You still don't get it. The goal is not an artificial creation to fill up the void left by the ``nihilism" of modern society. Weightlifting won't do it, or else football teams would be centers of enlightenment. Moreover, you are still thinking in categories that we ourselves have rejected, so an ``honest" conversation with you is still not possible.

It is not for nothing that Evola titled one of his books: ``The Myth of Blood".


\hfill

\texttt{Lionel on 2017-12-07 at 12:57 said: }

Putting together a few threads here.

My brother, Bhante Jason, is an Early Buddhist monk. He eats once a day, is homeless, spends Rains in a cave, owns only his bowl and robes, walks around everywhere barefoot, etc. He ordained in a highly respected Sri Lankan monastery, saw the failings there first hand, and walked off, finding a way back to Australia to instantiate the old way is not just possible without cultural support, but spiritually, mentally and physically superior. He never heard of Evola before my introduction.

Original Buddhism is non Traditional as Guenon identified, because Buddha denied the need for succession and lineage (apart from the suttas themselves) and stated clearly that there are no secret teachings. But it is also very supportive of Tradition in that he encouraged Brahmins to continue being Brahmins, showing a way his crystalised ascesis could make one the best Brahmin possible in this cycle. He also spoke of how vital it is to respect ones own ancestors, and only go out to the Homeless life with the blessing of ones parents and family. 

A difficult road that not many can walk, Early Buddhism is the Anarch that redeems the necessary castes by transcending them. What he has taught me thus turned this family man into quite the devout, soon to be Baptised Anglican.


\hfill

\texttt{Cologero on 2022-02-06 at 19:27 said: }

Lionel, Ananda Coomaraswamy seems to think otherwise, in Hinduism and Buddhism\footnote{\url{https://www.gornahoor.net/library/hinduism_buddhism.pdf}} where he writes:

\begin{quotex}
The more superficially one studies Buddhism, the more it seems to differ from the Brahmanism in which it originated; the more profound our study, the more difficult it becomes to distinguish Buddhism from Brahmanism, or to say in what respects, if any, Buddhism is really unorthodox.
\end{quotex}


\hfill

\texttt{Lionel Chan on 2022-02-06 at 20:06 said: }

Peace Sir,

Thank you for your engagement. Much has changed in my brother's life and my own since 2017, let alone the state of the world… the reminder of the strange twists in the path is very welcome.

I don't see My Coomaswarmy's comment as antithetical to what I wrote about what WAS my brother's Dao at that time, and still may be in some form. In times of degeneration, a kind of ``Protestant" and Fundamentalist" call can be very Traditional indeed. And as Mr Guenon emphasises, certainly it is rarer for the ``spiritual genius" to reach Deliverance with no initiation within an orthodoxy (a Left Hand way perhaps), but it is not excluded from possibility.

Salaam.


\end{sffamily}\end{footnotesize}
